% X_0^3390(1)/<w_5,w_6,w_113> is hyperelliptic: a genus-four double cover of the genus-two star, branched at the
% rational CM points of discriminants -3 and -40, which form a fibre of the star's hyperelliptic map.  The fibre
% statement is a dual test (paper/hyperelliptic-15-38.tex, Section 2) on a star whose Jacobian has no rational torsion.
% Tests: tests/Hyperelliptic3390_1.m and tests/_offline/Hyperelliptic3390_1.m.
\documentclass[11pt,reqno]{amsart}
\usepackage[margin=1.15in]{geometry}
\usepackage{amsmath,amssymb,amsthm,mathtools}
\usepackage[colorlinks=true,linkcolor=blue,citecolor=blue,urlcolor=blue]{hyperref}
\theoremstyle{plain}
\newtheorem{theorem}{Theorem}[section]
\newtheorem{lemma}[theorem]{Lemma}
\theoremstyle{definition}
\newtheorem{remark}[theorem]{Remark}
\DeclareMathOperator{\Jac}{Jac}
\DeclareMathOperator{\Pic}{Pic}
\newcommand{\Q}{\mathbb{Q}}
\newcommand{\Z}{\mathbb{Z}}
\newcommand{\F}{\mathbb{F}}
\begin{document}
\title{$X_0^{3390}(1)/\langle w_5,w_6,w_{113}\rangle$ is hyperelliptic}
\author{}
\date{\today}
\maketitle

\section{Statement}

Let $X=X_0^{3390}(1)$, $W$ its Atkin--Lehner group of order $16$, $H=\langle w_5,w_6,w_{113}\rangle$, the only
index-two subgroup for which $Y=X/H$ has genus four, and $X^*=X/W$ the star, of genus two. The double cover
$Y\to X^*$ has two branch points by Riemann--Hurwitz; a point of $X^*$ is a branch point exactly when its orbit
on $X$ consists of fixed points of some $w_m$ with $m\notin H$, and Ogg's count of those fixed points by CM
discriminant finds them only for $w_3$ (discriminant $-3$) and $w_{10}$ (discriminant $-40$), eight each. So the
branch points are the rational CM points $P_{-3}$ and $P_{-40}$ of $X^*$.

\begin{theorem}\label{thm:main}
$Y$ is hyperelliptic.
\end{theorem}

By the normal form of \texttt{paper/star-2002-proof.tex}, Lemma~8, a genus-four double cover of a genus-two curve
branched at two points is hyperelliptic if and only if the two branch points form a fibre of the hyperelliptic map
of the base. So the theorem says that $P_{-3}+P_{-40}$ is a fibre of the $g^1_2$ of $X^*$, i.e.\ that
$P_{-3}+P_{-40}-F$ is principal for one, hence every, fibre $F$.

\section{Proof}

\emph{The star has no rational torsion.} The orders $\#\Jac(X^*)(\F_p)=L_p(1)$, from the point counts of $X^*$
over $\F_p$ and $\F_{p^2}$ by the trace formula, have greatest common divisor $1$ over the good primes below
$120$; so $\Jac(X^*)(\Q)_{\mathrm{tors}}=0$ and a degree-zero divisor on $X^*$ is principal as soon as it is
torsion.

\emph{The pairing is torsion.} The dual test of \texttt{paper/hyperelliptic-15-38.tex}, Proposition~2.1, at
$M=4\cdot1695=6780$ (\texttt{tests/\_offline/Hyperelliptic3390\_1.m}): let $W$ be the space of cusp forms of
weight $\tfrac32$ on $\Gamma_0(6780)$ vanishing at every index up to the precision that carries no CM point. It
contains the two forms $\Phi_\varphi$ attached to the star, and at precision $1100$ it has dimension $51$, too
large for the Hecke operator $T(49)$ to be read off the coefficients available (its characteristic polynomial on
the star's two forms is $x^2+7x+11$, the polynomial of the traces of Frobenius of $X^*$ at $7$). The test
therefore uses the whole of $W$, which is the sufficient direction of the dual test: a principal part with
divisor $P_{-3}+P_{-40}-P_{-43}-P_{-67}$ in the dictionary of Guo--Yang pairs to zero with every form of $W$, as
does one with divisor $P_{-3}+P_{-40}-P_{-123}-P_{-168}$ (and, as a consistency check,
$P_{-43}+P_{-67}-P_{-123}-P_{-168}$), while $P_{-3}-P_{-40}$ and $P_{-3}+P_{-43}-P_{-40}-P_{-67}$ are refused. The point $P_{-3}$ is elliptic of order three on $X$ and fixed
by $w_3$; the dictionary carries the weight $1/e_d$ with $e_{-3}=6$ on $X^*$, which is why the test is run in
this form rather than through a Borcherds form of weight two.

So $P_{-3}+P_{-40}-P_{-43}-P_{-67}$ is torsion, hence principal: $P_{-3}+P_{-40}$ and $P_{-43}+P_{-67}$ are
fibres of one map of degree two, which is the hyperelliptic map of $X^*$. By the normal form, $Y$ is
hyperelliptic. \qed

\begin{remark}
The fibres $\{-43,-67\}$ and $\{-123,-168\}$ were found numerically, by the Shimizu screen of the invariant
differentials of $X^*$ at its rational CM points; the proof uses them only as the other side of a relation that
the dual test certifies, so nothing depends on the numerics. Since $\Jac(X^*)$ has rank two, equal to its
dimension, the rational points of $X^*$ cannot be listed by Chabauty's method, and the fibre statement is not
accessible that way.
\end{remark}

\end{document}
